\chapter{Hostel Care}

\section*{Definition and Overview}
Hostel care, now known in Australia as residential aged care at a lower level of support, provides accommodation, meals, personal care, and some nursing support for older people who can no longer manage independently at home but do not need full nursing care. It is one of the options considered when a person's care needs exceed what can be provided at home. In modern Australian aged care, the former distinction between hostels and nursing homes has been replaced by a single residential aged care system with services tailored to each resident's assessed needs. Decisions about entering residential care are significant and are best made with full information and family involvement.

\section*{Epidemiology}
Around 200,000 older Australians live in residential aged care at any time, and about 10\% of people over 70 will spend time in residential care. The average age at entry is around 85. Most residents need help with personal care, medication, and daily activities, and many have dementia. The population of older Australians is growing rapidly, increasing demand for aged care, while most older people continue to prefer and remain living at home with support for as long as possible.

\section*{Aetiology and Pathophysiology}
\begin{itemize}
    \item \textbf{Frailties of ageing:} declining physical and cognitive function increases care needs
    \item \textbf{Common reasons for entry:} dementia, falls, incontinence, inability to manage medications, and carer exhaustion
    \item \textbf{Assessment:} entry follows a formal aged care assessment of needs
    \item \textbf{Services provided:} accommodation, meals, personal care, laundry, activities, and nursing support
    \item \textbf{Levels of care:} services are tailored, with higher-level nursing care for greater needs
    \item \textbf{Pathophysiology of need:} loss of independence in activities of daily living makes safe living at home impossible
\end{itemize}

\section*{Clinical Features}
\begin{itemize}
    \item Difficulty with daily tasks: bathing, dressing, toileting, and eating
    \item Unsafe mobility, frequent falls, and wandering
    \item Medication mismanagement and poor adherence
    \item Incontinence and pressure injury risk
    \item Social isolation and loneliness at home
    \item Carer burnout and the inability of family to sustain care
    \item Cognitive decline requiring supervision
\end{itemize}

\section*{Differential Diagnosis}
\begin{itemize}
    \item The decision between staying at home with support, residential care, and higher-level nursing care depends on assessed needs
    \item Some people need only short-term respite care rather than permanent placement
    \item Home care packages meet the needs of many people who might otherwise enter residential care
    \item The choice of facility is guided by location, services, and availability
    \item Assessment and family discussion clarify the most appropriate option
\end{itemize}

\section*{When to Seek Medical Care}
Contact My Aged Care to arrange an aged care assessment, and speak with a GP about care needs. Seek urgent care for falls, acute illness, or any sudden decline.

\textbf{Call 000 (or local emergency services) for:}
\begin{itemize}
    \item A fall with a suspected fracture or head injury
    \item Collapse, severe confusion, or reduced consciousness
    \item Chest pain, severe breathlessness, or stroke symptoms
    \item Signs of severe infection with high fever
    \item A carer in crisis or any emergency situation
\end{itemize}

\section*{Diagnosis and Investigations}
\begin{itemize}
    \item \textbf{Aged Care Assessment:} a comprehensive assessment by an Aged Care Assessment Team determines eligibility for residential care
    \item \textbf{Medical review:} the GP assesses physical, cognitive, and functional status
    \item \textbf{Needs assessment:} care needs are matched to services
    \item \textbf{Facility tour:} families visit and compare facilities
    \item \textbf{Financial assessment:} fees and means testing are explained
    \item \textbf{Review:} care plans are reviewed regularly after entry
\end{itemize}

\section*{Management}
\begin{itemize}
    \item \textbf{Choosing a facility:} location, quality, services, and availability are considered
    \item \textbf{Care planning:} an individualised plan covers personal care, medication, activities, and health
    \item \textbf{Medication management:} staff administer and review medicines
    \item \textbf{Nursing and allied health:} nursing care, physiotherapy, and other services are provided as needed
    \item \textbf{Activities and social programs:} promote engagement and wellbeing
    \item \textbf{Family involvement:} regular contact and family meetings
    \item \textbf{Transition support:} adjustment to the new environment is supported
    \item \textbf{Quality and complaints:} residents and families can raise concerns through the Aged Care Quality and Safety Commission
\end{itemize}

\section*{Complications}
\begin{itemize}
    \item Falls, fractures, and pressure injuries
    \item Infections, including chest and urinary infections, and COVID-19 outbreaks
    \item Medication errors and adverse effects
    \item Worsening of dementia and behavioural symptoms
    \item Depression, anxiety, and loneliness
    \item Malnutrition and dehydration
    \item Adjustment difficulties for residents and families
\end{itemize}

\section*{Prognosis and Outcomes}
Residential care provides a safe environment for people who cannot manage at home, with 24-hour support and access to nursing care. Quality of life varies with the facility and the resident's condition, but good care, meaningful activities, and family involvement are associated with better wellbeing. Many residents remain stable for years, while those with progressive conditions receive increasing support. The quality of care and the relationship between staff, residents, and families are the strongest predictors of good outcomes.

\section*{Prevention and Self-Care}
\begin{itemize}
    \item Plan ahead: discuss future care preferences with family before a crisis
    \item Contact My Aged Care early for assessment, as waiting lists exist
    \item Visit and compare facilities before choosing
    \item Involve the older person in decisions as much as possible
    \item Maintain health: immunisations, nutrition, and physical activity
    \item Stay connected with family and community
    \item Understand fees, rights, and the complaints process
    \item Reassess needs regularly and adjust care
\end{itemize}

\section*{Sources and Further Reading}
The following reputable sources provide further information for healthcare students and patients seeking reliable, current advice about residential aged care.

\begin{itemize}
    \item My Aged Care. \textit{Residential aged care and entry process.} https://www.myagedcare.gov.au/
    \item Aged Care Quality and Safety Commission. \textit{Quality standards and complaints.} https://www.agedcarequality.gov.au/
    \item Older Persons Advocacy Network (OPAN). \textit{Advocacy and support for aged care residents.} https://opan.com.au/
    \item Healthdirect. \textit{Aged care services and support.} https://www.healthdirect.gov.au/
    \item Australian Government Department of Health. \textit{Aged care reform information.} https://www.health.gov.au/
    \item National Ageing Research Institute. \textit{Evidence on aged care and ageing well.} https://www.nari.net.au/
\end{itemize}
